% Chapter Template

\chapter{Conclusiones} % Main chapter title

\label{Chapter5}
En el presente capítulo se exponen las conclusiones más relevantes del trabajo realizado junto con los próximos pasos a seguir.

%----------------------------------------------------------------------------------------

%----------------------------------------------------------------------------------------
%	SECTION 1
%----------------------------------------------------------------------------------------

\section{Conclusiones generales}

El trabajo permitió obtener un prototipo funcional de un router MIDI de cuatro entradas y ocho salidas, con ruteo y filtrado configurables. El desarrollo incluyó el firmware embebido, la aplicación de escritorio para configuración, las placas electrónicas y el montaje mecánico en un gabinete. Como resultado, se obtuvo un sistema completo y operativo.

Los principales resultados del trabajo pueden resumirse en los siguientes puntos:
\begin{itemize}
	\item Se cumplieron los requerimientos funcionales principales planteados para el sistema. En particular, se implementó la cantidad prevista de entradas y salidas, el almacenamiento de presets y la configuración del comportamiento del equipo mediante ruteo y filtrado de mensajes.

	\item El prototipo mostró un funcionamiento satisfactorio en la práctica. Los ensayos realizados permitieron verificar la operación general del sistema y validar su desempeño en distintos escenarios de uso.

	\item Desde el punto de vista temporal, el equipo presentó una respuesta adecuada para la aplicación prevista. En todos los casos analizados, la latencia permaneció por debajo del criterio de aceptación adoptado, incluso en el escenario más exigente.

	\item Algunos aspectos previstos en la planificación original no se completaron en esta etapa y quedaron planteados como líneas de trabajo futuro. Asimismo, el desarrollo demandó un plazo de ejecución mayor al previsto originalmente. No obstante, estas cuestiones no invalidan el resultado principal alcanzado ni el cumplimiento de los objetivos técnicos centrales del proyecto.
\end{itemize}

%----------------------------------------------------------------------------------------
%	SECTION 2
%----------------------------------------------------------------------------------------

\section{Próximos pasos}

Para avanzar desde el prototipo actual hacia una versión más madura y comercializable del sistema, se identificaron los siguientes puntos de mejora:
\begin{itemize}
	\item Incorporar la exposición de las entradas y salidas MIDI a través del puerto USB, para habilitar el control desde la computadora y hacia la computadora.
	\item Automatizar el proceso de compilación y testing a través de un pipeline de CI/CD (\textit{Continuous Integration/Continuous Deployment}), para poder ofrecer un mecanismo de versionado y trazabilidad más estructurado.
	\item Reducir copias de memoria en la propagación de mensajes, ya que el despacho actual utiliza múltiples colas con duplicación de bytes.
	\item Dado que la alimentación es por el puerto USB, incorporar un mecanismo de almacenamiento de energía y detección de desconexión para guardar cambios no aplicados y restablecer la configuración cuando se reinicie el equipo.
	\item Evaluar el uso de una variante de 100 pines del microcontrolador, con el objetivo de reducir costo y área de circuito impreso.
	\item Eliminar componentes prescindibles para la versión final del producto.
	\item Refinar la experiencia de uso de la aplicación de escritorio, y agregar la posibilidad de actualizar el firmware desde la misma.
	\item Redactar un manual de usuario.
	\item Diseñar la serigrafía del frente del equipo.
\end{itemize}
